\documentclass[10pt, a4paper]{article}

% --- Packages ---
\usepackage[utf8]{inputenc}
\usepackage[T1]{fontenc}
\usepackage[margin=1in]{geometry}
\usepackage{xcolor}
\usepackage{enumitem}
\usepackage{titlesec}
\usepackage{hyperref}
\usepackage{charter}
\usepackage{parskip}

% --- Colors ---
\definecolor{primarycolor}{RGB}{0, 51, 102}
\hypersetup{colorlinks=true, urlcolor=primarycolor, linkcolor=primarycolor}

% --- Layout ---
\pagestyle{empty}

\begin{document}

% --- Sender Information ---
\begin{flushright}
    \textbf{Kartavya Niraj Dikshit} \\
    Halbmondstraße 7 \\
    91054 Erlangen \\
    (+49) 155 10940829 \\
    \href{mailto:kartavya.dikshit@gmail.com}{kartavya.dikshit@gmail.com} \\
    \today
\end{flushright}

\vspace{1em}

% --- Recipient Information ---
\textbf{Siemens AG} \\
Smart Infrastructure \\
Recruiting Team \\
Nürnberg, Deutschland

\vspace{2em}

% --- Subject Line ---
{\large \textbf{\color{primarycolor} Bewerbung als Werkstudent (f/m/d) Data Scientist}} \\
\textbf{Job ID: 498945}

\vspace{1.5em}

Sehr geehrte Damen und Herren,

als Masterstudent der Data Science an der Friedrich-Alexander-Universität Erlangen-Nürnberg mit einer Spezialisierung auf künstliche Intelligenz und Machine Learning verfolge ich die Entwicklungen von Siemens im Bereich Smart Infrastructure mit großer Begeisterung. Ihre Vision, durch Daten und State-of-the-Art-Modelle reale industrielle Herausforderungen zu lösen, deckt sich exakt mit meinem akademischen Fokus und meiner bisherigen technischen Erfahrung.

In meinem aktuellen Studium an der FAU vertiefe ich mein Wissen in den Bereichen Deep Learning und Scalable Data Architectures. Besonders der Fokus auf „Industrial AI“ ermöglicht es mir, theoretische Modelle in den Kontext komplexer Infrastrukturen zu setzen. Diese akademische Grundlage möchte ich nutzen, um Ihr Team beim Benchmarking von Modellen und der kontinuierlichen Verbesserung Ihres Modell-Repositorys tatkräftig zu unterstützen.

Durch meine bisherigen Tätigkeiten als AI/ML Engineer Intern verfüge ich über fundierte Kenntnisse in der Entwicklung von Prototypen und deren Überführung in produktionsreife Umgebungen. Ich beherrsche nicht nur die gängigen ML-Frameworks wie PyTorch und Scikit-learn, sondern lege auch großen Wert auf Software-Engineering-Prinzipien wie Clean Code, OOP und eine saubere Versionsverwaltung via Git. In meinen Projekten konnte ich bereits zeigen, dass ich komplexe Ergebnisse durch aussagekräftige Visualisierungen transparent kommunizieren kann – eine Fähigkeit, die für die Wissensvermittlung innerhalb Ihres Data Science Teams essenziell ist.

Siemens steht für Innovation und Qualität. Ich bin motiviert, meine Eigeninitiative und meinen Growth Mindset in ein Team einzubringen, das die Infrastruktur von morgen gestaltet. Die Möglichkeit, 14 bis 20 Stunden pro Woche an zukunftsweisenden Projekten in Nürnberg mitzuwirken, sehe ich als idealen nächsten Schritt in meiner beruflichen Entwicklung.

Gerne überzeuge ich Sie in einem persönlichen Gespräch von meiner Eignung und meiner Begeisterung für diese Position.

Mit freundlichen Grüßen

\vspace{1em}

Kartavya Niraj Dikshit

\end{document}